\section{USAJMO 2018}
        \begin{enumerate}
        	\item (\textit{USAJMO 2018}) Note: For any geometry problem whose statement begins with an asterisk ($*$), the first page of the solution must be a large, in-scale, clearly labeled diagram. Failure to meet this requirement will result in an automatic 1-point deduction. 


 For each positive integer $n$, find the number of $n$-digit positive integers that satisfy both of the following conditions:


 $\bullet$ no two consecutive digits are equal, and


 $\bullet$ the last digit is a prime.


 

	\item (\textit{USAJMO 2018}) Note: For any geometry problem whose statement begins with an asterisk ($*$), the first page of the solution must be a large, in-scale, clearly labeled diagram. Failure to meet this requirement will result in an automatic 1-point deduction. 


 Let $a,b,c$ be positive real numbers such that $a+b+c=4\sqrt[3]{abc}$. Prove that 
 \[2(ab+bc+ca)+4\min(a^2,b^2,c^2)\ge a^2+b^2+c^2.\]


 

	\item (\textit{USAJMO 2018}) Note: For any geometry problem whose statement begins with an asterisk ($*$), the first page of the solution must be a large, in-scale, clearly labeled diagram. Failure to meet this requirement will result in an automatic 1-point deduction. 


 ($*$) Let $ABCD$ be a quadrilateral inscribed in circle $\omega$ with $\overline{AC} \perp \overline{BD}$. Let $E$ and $F$ be the reflections of $D$ over lines $BA$ and $BC$, respectively, and let $P$ be the intersection of lines $BD$ and $EF$. Suppose that the circumcircle of $\triangle EPD$ meets $\omega$ at $D$ and $Q$, and the circumcircle of $\triangle FPD$ meets $\omega$ at $D$ and $R$. Show that $EQ = FR$.


 

	\item (\textit{USAJMO 2018}) Note: For any geometry problem whose statement begins with an asterisk ($*$), the first page of the solution must be a large, in-scale, clearly labeled diagram. Failure to meet this requirement will result in an automatic 1-point deduction. 


 Triangle $ABC$ is inscribed in a circle of radius 2 with $\angle ABC \geq 90^\circ$, and $x$ is a real number satisfying the equation $x^4 + ax^3 + bx^2 + cx + 1 = 0$, where $a=BC,b=CA,c=AB$. Find all possible values of $x$.


 

	\item (\textit{USAJMO 2018}) Note: For any geometry problem whose statement begins with an asterisk ($*$), the first page of the solution must be a large, in-scale, clearly labeled diagram. Failure to meet this requirement will result in an automatic 1-point deduction. 


 Let $p$ be a prime, and let $a_1, \dots, a_p$ be integers. Show that there exists an integer $k$ such that the numbers

 \[a_1 + k, a_2 + 2k, \dots, a_p + pk\]produce at least $\frac{1}{2} p$ distinct remainders upon division by $p$.


 

	\item (\textit{USAJMO 2018}) Note: For any geometry problem whose statement begins with an asterisk ($*$), the first page of the solution must be a large, in-scale, clearly labeled diagram. Failure to meet this requirement will result in an automatic 1-point deduction. 


 Karl starts with $n$ cards labeled $1,2,3,\dots,n$ lined up in a random order on his desk. He calls a pair $(a,b)$ of these cards swapped if $a>b$ and the card labeled $a$ is to the left of the card labeled $b$. For instance, in the sequence of cards $3,1,4,2$, there are three swapped pairs of cards, $(3,1)$, $(3,2)$, and $(4,2)$.


 He picks up the card labeled 1 and inserts it back into the sequence in the opposite position: if the card labeled 1 had $i$ card to its left, then it now has $i$ cards to its right. He then picks up the card labeled $2$ and reinserts it in the same manner, and so on until he has picked up and put back each of the cards $1,2,\dots,n$ exactly once in that order. (For example, the process starting at $3,1,4,2$ would be $3,1,4,2\to 3,4,1,2\to 2,3,4,1\to 2,4,3,1\to 2,3,4,1$.)


 Show that, no matter what lineup of cards Karl started with, his final lineup has the same number of swapped pairs as the starting lineup.


 

\end{enumerate}